\begin{lemma}[Finite-scale calibration]\label{lem:finite-calibration}
Let $a=2d\kappa$, $b=\nabla\rho$, and let $D\subset B(0,R)$ be measurable,
$R\geq1$. Put $u=U_D$ and suppose
\[
 u\geq-\delta\quad\hbox{on }B(0,2R),\qquad
 |u|\leq\delta\quad\hbox{on }B(0,2R)\setminus D.
\]
Then $w=(u-\delta)_+$ vanishes off $D$, belongs to $H^1(\mathbb R^d)$,
and, with $A=\{u>\delta\}$,
\begin{equation}\label{eq:finite-calibration}
 \int_{\mathbb R^d}|\nabla w+a\mathbf1_A b|^2
 +a^2\int_{D\setminus A}|b|^2
 \leq C\delta R^{d-1}.
\end{equation}
In particular $|D\setminus A|\leq C\delta R^{d-1}$.
\end{lemma}
\begin{proof}
The potential is continuous, tends to zero at infinity, and is harmonic
off $\overline D$. The maximum principle extends $|u|\leq\delta$ from the
annulus to the exterior of $B(0,R)$. The potential equation is
$\operatorname{div}(\nabla u+a\mathbf1_D b)=0$.
Since $\mathbf1_D b\in L^2$, the usual Fourier multiplier estimate gives
$\nabla u\in L^2$; hence the compactly supported truncation $w$ is in $H^1$.
Testing the equation with $w$, using $A\subset D$, gives
\begin{equation}\label{eq:calibration-test-w}
 \int|\nabla w|^2=-a\int b\cdot\nabla w.
\end{equation}
Choose a smooth $\chi$ equal to one on $B(0,R)$, zero off $B(0,2R)$,
with $|\nabla\chi|\leq C/R$ and $|\Delta\chi|\leq C/R^2$.
Put $v=u-w$, so $|v|\leq\delta$. The equation tested with $\chi\rho$,
together with \eqref{eq:calibration-test-w}, gives the exact identity
\[
 \int|\nabla w+a\mathbf1_D b|^2
 =a\int b\cdot\nabla w+a^2\int_D|b|^2
 =-a\int\nabla v\cdot\nabla(\chi\rho)
 =a\int v\,d\Delta(\chi\rho).
\]
The convexity measure $\mu=\Delta\rho$ has
$\mu(B(0,2R))\leq CR^{d-1}$. Since
$\Delta(\chi\rho)=\chi\mu+2b\cdot\nabla\chi+\rho\Delta\chi$,
its total variation is at most $CR^{d-1}$. Integration by parts is
justified by uniform and local $H^1$ mollification of the continuous
function $v$. Bounding the right side by $C\delta R^{d-1}$ and splitting
the square over $A$ and $D\setminus A$ proves the estimate.
Finally $|b|\geq c$ almost everywhere, by Euler's identity and
$\rho(x)\geq c|x|$.
\end{proof}

\begin{corollary}[Exact rigidity]\label{cor:calibration-rigidity}
If $U_D\geq0$ and $U_D=0$ on $D^c$, then, up to null sets, $D=rK$ for
some $r\geq0$ and $U_D=2d\kappa(r-\rho)_+$.
\end{corollary}
\begin{proof}
If $|D|=0$, take $r=0$. Otherwise take $\delta=0$. Then $D=\{u>0\}$ up to null sets and
$\nabla u=-a\mathbf1_{\{u>0\}}\nabla\rho$.
Thus $u$ is Lipschitz and $u+a\rho$ is constant on each connected
component of its positivity set. On such a component write
$u=a(r-\rho)$. The segment from any point of the component toward the
origin cannot leave it: at a first exit its value would still be
strictly positive. Every component therefore contains the origin, so
there is only one. Its boundary lies in $\{\rho=r\}$; following each
radial segment now shows that the component is $rK$.
\end{proof}

\begin{corollary}[Few sites have small local time]\label{cor:low-occupation}
On the coarse event, let $\delta_n$ be a uniform bound for
$|\widetilde\ell_n-U_{D_n}|$. Then simultaneously for every $t\geq0$,
\[
 \#\{x\in A_n:\ell_n(x)\leq t\}
 \leq C(t+\delta_n)r_n^{d-1}.
\]
\end{corollary}
\begin{proof}
Apply the lemma with truncation threshold $t+\delta_n$ and $R=Cr_n$.
Every cell with $\ell_n(x)\leq t$ lies in $\{U_{D_n}\leq t+\delta_n\}$.
\end{proof}
